% GPT-6 Astra Pro; /root/fixed_surface; 22 November 2026 UTC.
% Effort not exposed. New complete proof for independent review.
\section[HD-F1-S24: reflexive local Ext and the actual normal obstruction]{HD-F1-S24: the actual self-Ext sheaves are reflexive exterior powers, and the residual obstruction lies in normal cohomology}
\label{sec:f1-norm-fourfold-reflexive-ext}

Retain the actual norm image $i:Y\hookrightarrow P$, its ideal $I$,
and $E=i_*\mathcal O_Y$ on the open family of S13--S22. Write
\[
 \mathscr N=\mathcal Hom_P(I,E)
       =\mathcal Hom_Y(I/I^2,\mathcal O_Y),\qquad
 \mathscr E^q=\mathcal Ext_P^q(E,E),
 \tag{F1S24.1}
\]
viewed as coherent sheaves on $Y$. The actual Yoneda maps give
\[
 \mathscr E^q\simeq\bigl(\textstyle\bigwedge^q\mathscr N\bigr)^{**}
       \quad(0\le q\le4),\qquad
 \mathscr E^q=0\quad(q>4).
 \tag{F1S24.2}
\]
In particular all five sheaves are reflexive, with respective generic
ranks $1,4,6,4,1$. This is an assertion about the actual completed
local rings and the actual coherent sheaves, not only their tangent
cones.

The entire nonscalar degree-two source of S23 has the exact sequence
\[
 0\longrightarrow H^1(Y,\mathscr N)
 \xrightarrow{j}\operatorname{Ext}^1_P(I,E)
 \xrightarrow{\ell}H^0\bigl(Y,(\bigwedge^2\mathscr N)^{**}\bigr)
 \xrightarrow{d_2}H^2(Y,\mathscr N).
 \tag{F1S24.3}
\]
The differential and all maps will be specified below. On the S23
open where the period derivative has rank twelve, its actual
four-dimensional residual map factors further as
\[
 (\ker\mu_C)^*\xrightarrow{\overline\eta_Y}H^1(Y,\mathscr N)
    \xrightarrow{j}\operatorname{Ext}^1_P(I,E),\qquad
                 j\overline\eta_Y=\overline h_E.
 \tag{F1S24.4}
\]
The first arrow is the quotient of the explicitly defined derivative
map in paragraph5. Its rank remains uncomputed. Neither the
reflexivity statement nor (F1S24.3) asserts a splitting of the global
derived self-intersection or injectivity of the full trace.

\paragraph{1. An actual smooth formal slice and its lifted resolution.}
At a smooth point of $Y$ the assertions are the usual Koszul
calculation for its actual regular embedding of codimension four.
Consider a singular point $y$, which is the image of a complete
pencil by S16--S17. Its embedding dimension is six. In the completed
regular local ring of $P$ choose two elements $z_1,z_2$ of its ideal
whose independent linear terms span the kernel of the cotangent
map to $Y$. Extend them to formal coordinates. Then
\[
 R_8=\widehat{\mathcal O}_{P,y}
      =\mathbf C[[x_1,\ldots,x_6,z_1,z_2]],\qquad
 S=\widehat{\mathcal O}_{Y,y}=R/J,
 \quad R=\mathbf C[[x_1,\ldots,x_6]],
 \tag{F1S24.5}
\]
and the original completed ideal is $(z_1,z_2)+JR_8$.
This follows by reducing its elements modulo the two contained
parameters; it is the given ambient quotient map.

Put $A=\mathbf C[x_1,\ldots,x_6]$. By S17 the actual associated
graded ring of $S$ is $A/J_0$, where $J_0$ is generated by the
three signed minors $q_1,q_2,q_3$ of one of the following matrices,
with the indicated unused polynomial coordinates:
\[
 \begin{aligned}
 (1,1,1):\quad B_0&=\begin{pmatrix}a&b&c\\d&e&f\end{pmatrix},\\
 (0,1,2):\quad B_0&=\begin{pmatrix}x_0&y_0&y_1\\x_1&y_1&y_2\end{pmatrix}
                  &&\text{one unused coordinate},\\
 (0,0,3):\quad B_0&=\begin{pmatrix}y_0&y_1&y_2\\y_1&y_2&y_3\end{pmatrix}
                  &&\text{two unused coordinates}.
 \end{aligned}
 \tag{F1S24.6}
\]
For a matrix with rows $(a,b,c),(d,e,f)$ use
$q=(bf-ce,cd-af,ae-bd)$, so $q B_0^t=0$.
The graded resolution is exactly
\[
 0\longrightarrow A(-3)^2\xrightarrow{B_0^t}A(-2)^3
       \xrightarrow{q}A\longrightarrow A/J_0\longrightarrow0.
 \tag{F1S24.7}
\]
Here is a direct verification common to all three cases.
The first map is injective because one of its two-by-two minors
is nonzero in the polynomial domain. Its cokernel maps onto
$J_0$. The Hilbert series of that cokernel is
$(3t^2-2t^3)/(1-t)^6$, while S17 gives
\[
 \operatorname{Hilb}(A/J_0,t)=\frac{1+2t}{(1-t)^4}
                       =\frac{1-3t^2+2t^3}{(1-t)^6}.
\]
Consequently $J_0$ has the same series as the cokernel, and the
surjection is an isomorphism degree by degree. This proves
(F1S24.7) using the actual presentations already obtained in S17.

Choose $f_i\in J$ with initial forms $q_i$.
The map $R^3\to J$ with row $f=(f_1,f_2,f_3)$ is strictly
surjective for the order filtrations, with its basis vectors
assigned order two. Indeed, for an element of $J$, express its
initial form by the $q_i$, subtract the corresponding combination
of $f_i$, and repeat in higher orders. Completeness gives convergent
formal coefficients and the asserted surjection at every filtration
level. Its kernel therefore has associated graded equal to the
kernel of $q$: an approximate graded syzygy is corrected by a
preimage of its higher-order error under this strict surjection.
Lift the two columns of $B_0^t$ to actual kernel elements. They
form a matrix $\Phi:R^2\to R^3$ whose linear part is $B_0^t$.
The same successive lifting shows that they generate the kernel.
They are independent, since a relation of least nonzero order
would give a relation under the injective first map of (F1S24.7).
Thus the actual complete local ring has the free resolution
\[
 0\longrightarrow R^2\xrightarrow{\Phi}R^3
     \xrightarrow{f}R\longrightarrow S\longrightarrow0,
       \qquad f\Phi=0.
 \tag{F1S24.8}
\]
All liftings are in the actual ideal and its actual syzygy module.
No identification of $S$ with its associated graded ring was used.

\paragraph{2. Computing the actual local self-Ext modules.}
Apply $\operatorname{Hom}_R(-,S)$ to (F1S24.8). Its first
differential is zero because $f_i=0$ in $S$. Hence
\[
 \operatorname{Ext}_R^0(S,S)=S,\quad
 \operatorname{Ext}_R^1(S,S)=M:=\ker(\Phi^t:S^3\to S^2),\quad
 \operatorname{Ext}_R^2(S,S)=\operatorname{coker}(\Phi^t:S^3\to S^2).
 \tag{F1S24.9}
\]
The last cokernel is the canonical module $\omega_S$ after choosing
the coordinate volume form of $R$. More precisely, finite-map
duality gives the normalized dualizing complex
$R\operatorname{Hom}_R(S,R[6])$. The ring $S$ is Cohen--Macaulay
of dimension four, by S16--S17, so this complex is
$\omega_S[4]$ and
\[
 \omega_S=\operatorname{Ext}_R^2(S,R)
          =\operatorname{coker}(\Phi^t:R^3\to R^2).
 \tag{F1S24.10}
\]
This module is annihilated by $J$: it is the cohomology of the
displayed dualizing complex as an $S$-complex. Reducing that
presentation modulo $J$ therefore gives exactly the last cokernel
in (F1S24.9), establishing the claimed isomorphism.

Extend (F1S24.8) to $R_8$, and tensor it with the Koszul resolution
of $z_1,z_2$. It resolves the actual $S$ over $R_8$: those two
parameters form a regular sequence on $S[[z_1,z_2]]$.
After applying $\operatorname{Hom}_{R_8}(-,S)$ the two Koszul
differentials are zero. Taking cohomology therefore gives the
actual local identifications
\[
\begin{array}{c|c|c}
q&\operatorname{Ext}_{R_8}^q(S,S)&\text{generic rank over }S\\ \hline
0&S&1\\
1&M\oplus S^2&4\\
2&\omega_S\oplus M^2\oplus S&6\\
3&\omega_S^2\oplus M&4\\
4&\omega_S&1
\end{array}
\qquad\operatorname{Ext}_{R_8}^{q}(S,S)=0\quad(q>4).
 \tag{F1S24.11}
\]
The decompositions depend on the displayed formal coordinates and
free-resolution bases; they are not asserted to glue as direct sums.
The module $M$ is the actual normal module of $S\subset\operatorname{Spec}R$,
by the Hom sequence of its ideal.

The normal module $M$ is reflexive over the normal domain $S$:
it is the kernel of a map from the reflexive module $S^3$ to
the torsion-free module $S^2$. The canonical module is maximal
Cohen--Macaulay and has rank one. Localization gives its $S_2$
condition, and its associated primes are minimal; hence it is
torsion-free and therefore reflexive on this normal domain.
All modules in (F1S24.11) are consequently reflexive.

\paragraph{3. Descent and the exact reflexive exterior algebra.}
The coherent Ext sheaves commute with completion, as follows by
applying Hom to a finite local free resolution and using flatness
of completion. For any finite module, its double dual also commutes
with this flat base change. The kernel and cokernel of its actual
double-dual map vanish if their completions vanish. Faithful
flatness therefore descends the reflexivity in (F1S24.11) to the
actual local rings of $Y$. Along the smooth locus the Koszul
calculation gives the same statement. This proves global
reflexivity of every $\mathscr E^q$ in (F1S24.2).

Let $Y^{\rm reg}\xrightarrow{j_0}Y$ be the actual smooth open.
Normality makes its complement have codimension at least two.
On that open the actual ideal triangle identifies
$\mathscr E^1=\mathscr N$, and the Koszul Yoneda product is
the exterior product of the normal bundle. The degree-one
identification holds globally as well, by the sheaf Hom sequence
of $I\to\mathcal O_P\to E$: the map
$\mathcal Hom_P(E,E)\to\mathcal Hom_P(\mathcal O_P,E)$
is the identity of $\mathcal O_Y$, while restriction to $I$ is zero.

The graded commutator $\alpha\beta+\beta\alpha$ of two local
degree-one Yoneda classes vanishes on $Y^{\rm reg}$.
Since $\mathscr E^2$ is torsion-free, it
vanishes everywhere. In characteristic zero their squares also
vanish. Thus the actual Yoneda product factors through
$\bigwedge^q\mathscr N\to\mathscr E^q$.
Taking the unique reflexive extension of its isomorphism on
$Y^{\rm reg}$ gives exactly (F1S24.2); equivalently both sides
are the pushforward of the same exterior normal bundle from
$Y^{\rm reg}$. These maps preserve products, because their two
product composites agree on that open and the targets are
torsion-free. This establishes the claimed global reflexive
exterior algebra, without asserting formality of its underlying
derived object.

\paragraph{4. The precise global obstruction sequence.}
Put $\mathcal K=R\mathcal Hom_P(I,E)$, regarded on $Y$ via the
target module structure. Equivalently use
$R\mathcal Hom_Y(Li^*I,\mathcal O_Y)$ and the actual adjunction.
Its degree-zero cohomology is $\mathscr N$. The ideal triangle
gives $\mathcal H^1(\mathcal K)=\mathscr E^2$.
The actual first Postnikov truncation has triangle
\[
 \mathscr N\longrightarrow\tau_{\le1}\mathcal K
       \longrightarrow\mathscr E^2[-1]
       \xrightarrow{\xi}\mathscr N[1].
 \tag{F1S24.12}
\]
The higher truncation has cohomology in degrees at least two,
so it does not change hypercohomology in degree one. Taking
hypercohomology of (F1S24.12) gives (F1S24.3), with
$d_2=H^1(\xi)$ under its displayed shifted identifications.
Thus $\ell$ is the actual local-Ext edge map, and $j$ is the
injection of normal cohomology, not a chosen complement.
This is also the low-degree part of the local-to-global Ext
spectral sequence.

In particular $\mathscr E^2$ has no nonzero subsheaf supported
on the singular locus, or on any proper closed subset of $Y$.
This follows from its torsion-freeness and integrality of $Y$.
It excludes constructing a global self-Ext class by first choosing
a skyscraper subsheaf of $\mathscr E^2$ at a pencil singularity.
It does not eliminate global classes from $H^1(\mathscr N)$
or compatible sections in the kernel of $d_2$.

\paragraph{5. The actual derivative into normal cohomology.}
There is the following sheaf morphism on the same embedding:
\[
 \nu:T_P|_Y\longrightarrow\mathscr N,
       \qquad\nu(v)(f)=v(f)|_Y\quad(f\in I).
 \tag{F1S24.13}
\]
It factors through $I/I^2$ by the Leibniz rule. In the formal
coordinates of (F1S24.5)--(F1S24.8), its components in
$M\oplus S^2$ are precisely
\[
 (v_x,v_z)\longmapsto
     \left(\left(\sum_{r=1}^6v_r\frac{\partial f_i}{\partial x_r}
                                          \right)_{i=1}^3,
                    v_{z_1},v_{z_2}\right).
 \tag{F1S24.14}
\]
The first vector belongs to $M$, since differentiating $f\Phi=0$
and reducing modulo $J$ gives $\Phi^t(v(f_i))=0$.
Thus this formula retains the actual higher-order equations $f_i$,
not just their initial quadrics.

Let
\[
 \eta_Y:H^1(P,T_P)\longrightarrow H^1(Y,T_P|_Y)
                     \xrightarrow{H^1(\nu)}H^1(Y,\mathscr N)
 \tag{F1S24.15}
\]
be restriction followed by this derivative. The first arrow is
an isomorphism here: $T_P$ is trivial and S21 proves that the
actual $H^1(\mathcal O_P)\to H^1(\mathcal O_Y)$ is an
isomorphism. Its source has dimension sixty-four.

We now compare (F1S24.15) with the exact S23 Atiyah map, including
its convention. Identify the conormal of the diagonal by
$1\otimes f-f\otimes1\mapsto df$, with the first factor
the input and the second the output, as in the Huybrechts--Thomas
diagonal definition used in S23. The classical first-jet extension
of $E$ is
\[
 0\longrightarrow E\otimes\Omega_P^1
         \longrightarrow J^1(E)\longrightarrow E\longrightarrow0.
 \tag{F1S24.16}
\]
Locally $J^1(E)=(E\otimes_{\mathbf C}\mathcal O_P)/I_\Delta^2$.
Smoothness makes the diagonal conormal locally free, so this
extension is obtained by the exact diagonal tensor construction
even though $E$ need not be locally free. The quotient unit
$\mathcal O_P\to E$ lifts to $J^1(E)$ by $r\mapsto1_E\otimes r$.
For $f\in I$ its value lies in the first term and is exactly
$df|_Y$, since $f\otimes1_E=0$. Therefore, for the actual
boundary $\partial:E\to I[1]$ of the ideal triangle,
\[
 \operatorname{At}(E)=(d|_I)[1]\partial.
 \tag{F1S24.17}
\]
This is an equality of actual extension classes with the specified
diagonal sign, proved by the lift just written.

For $\kappa\in H^1(P,T_P)$, compose $d|_I$ with the
restriction of $\kappa:\Omega_P^1\to\mathcal O_P[1]$.
Its class in $\operatorname{Ext}^1_P(I,E)$ is
$j\eta_Y(\kappa)$. One explicit verification uses an affine
cover trivializing the ambient first-order deformation: if
$v_{ab}$ is its derivation cocycle, the resulting Hom cocycle
is $f\mapsto v_{ab}(f)|_Y$, exactly (F1S24.13).
Equivalently the same assertion follows by applying the
edge map to this degree-one Cech representative. Consequently
\[
 \operatorname{ob}_E(\kappa)
    =\bigl(j\eta_Y(\kappa)\bigr)[1]\partial
    =d^{\rm ideal}_2\,j\eta_Y(\kappa),
 \tag{F1S24.18}
\]
where $d^{\rm ideal}_2$ is the injective second-summand map
in S23, not the spectral-sequence differential in (F1S24.3).
In particular the local-Ext edge of this ideal obstruction is zero.

\paragraph{6. The exact remaining four-dimensional and full-trace gates.}
Both $j$ and $d^{\rm ideal}_2$ in (F1S24.18) are injective.
The actual norm family in S23 therefore proves
$\eta_Y\circ d\operatorname{Pr}=0$. Its exact period sequence
then gives the unique map $\overline\eta_Y$ in (F1S24.4),
with the defining equation
\[
 \overline\eta_Y\bigl(\jmath(\kappa)|_{\ker\mu_C}\bigr)
                       =\eta_Y(\kappa),\qquad\kappa\in\mathfrak D_P.
 \tag{F1S24.19}
\]
The equality is independent of the representative by the just-proved
vanishing on the actual period image. The full character identity
of S23 further gives
\[
          \operatorname{im}\overline\eta_Y
                     \subset\ker(\tau_Y\circ j).
 \tag{F1S24.20}
\]
The full degree-two trace kernel remains exactly $\ker\tau_Y$
on the actual middle group of (F1S24.3). That group is now
described by normal cohomology and the actual Postnikov-compatible
sections of the reflexive exterior normal sheaf. Neither its
dimension nor the ranks of $\tau_Y$, $d_2$ and the residual
map (F1S24.19) are assigned here.

\paragraph{Primary inputs and scope.}
The graded presentations and Hilbert series are the actual
S17 conclusions; their resolutions and the filtered lifting are
proved above. The precise finite-map duality used in
(F1S24.10) is Stacks Lemma47.15.8, Tag0AX0, whose complete proof
was read. Its Cohen--Macaulay concentration and canonical-module
property are Lemma47.20.2, Tag0AWS, with the full linked proof
of Lemma47.16.7, Tag0B5A. Reflexive kernels and extension from
the normal big open use Lemmas15.24.5 and15.24.18, Tags0EB8
and0AVB, with their complete proofs, and Lemma31.13.12.
The actual hypercohomology filtration is Stacks Section20.43,
Tag0BQP, and Example20.29.3, Tag0BKM; (F1S24.12) writes the
specific truncation that supplies the needed low-degree arrows.
The Atiyah convention and first-jet comparison are the diagonal
maps of Huybrechts--Thomas, corrected2013 version, Section2.2,
especially Lemma2.2, Definition2.3 and equation(2.7).
All these statements act on the original norm image and its
actual ideal. No simultaneous resolution of a deformed image
or exhaustion by deformations of covers is assumed.
